\documentclass[11pt]{article}
\usepackage[margin=1in]{geometry}
\usepackage{amsmath,amssymb,amsthm}
\usepackage[hidelinks]{hyperref}
\newcommand{\LL}{\mathcal L}
\title{A rational counterexample to infinite log-concavity\\
\large Complete solution of MathDB 347493}
\author{Principia Math in Progress}
\date{12 September 2026}
\begin{document}
\maketitle

\paragraph{Task 1. Resolve the infinite log-concavity conjecture.}
For $n\geq0$, put
\[
 f_n(z)=\sum_{j=0}^n\binom nj z^{\binom j2},
 \qquad
 \LL(a)_n=a_n^2-a_{n-1}a_{n+1}.
\]
Dilcher and Ulas conjectured that $(f_n(z))_{n\geq0}$ is infinitely
log-concave for every $0\leq z\leq1$.
This is false: at $z=3/10$, the sixth iterate satisfies
\begin{equation}\label{negative}
  \LL^6(f)_6(3/10)<-\frac{63}{10^{36}}<0.
\end{equation}
Only $f_0,\ldots,f_{12}$ enter this counterexample, so no convention
about negative indices is needed.

\paragraph{A finite exact certificate.}
Define rational numbers
\[
 a_{0,n}=\sum_{j=0}^n\binom nj(3/10)^{\binom j2}\quad(0\leq n\leq12),
 \qquad
 a_{r+1,n}=a_{r,n}^2-a_{r,n-1}a_{r,n+1}.
\]
At level $r$ we retain $r\leq n\leq12-r$. By induction,
$a_{r,n}=\LL^r(f)_n(3/10)$.
Exact rational calculation gives the following strict bounds:
\[
\begin{array}{c|ccc}
n&5&6&7\\ \hline
\text{lower bound for }10^{18}a_{5,n}&8&44&261\\
\text{upper bound for }10^{18}a_{5,n}&9&45&262
\end{array}
\]
In particular, the three entries are positive, and hence
\[
 a_{6,6}=a_{5,6}^2-a_{5,5}a_{5,7}
 <\frac{45^2-8\cdot261}{10^{36}}
 =-\frac{63}{10^{36}}.
\]
This proves \eqref{negative}.

For completeness, the exact calculation establishing every bound in
the table is reproduced below. This is a finite rational-arithmetic
certificate; the code uses only Python's standard library, with no
floating point operations.
\par\noindent\begin{minipage}{\linewidth}\small
\begin{verbatim}
from fractions import Fraction as F
from math import comb

q = F(3, 10)
a = [sum((F(comb(n,j))*q**comb(j,2)
          for j in range(n+1)), F(0))
     for n in range(13)]
for r in range(5):
    a = [a[i]**2-a[i-1]*a[i+1]
         for i in range(1, len(a)-1)]

for x, lo, hi in zip(a, (8,44,261), (9,45,262)):
    assert lo < 10**18*x < hi
assert a[1]**2-a[0]*a[2] < F(-63, 10**36)
\end{verbatim}
\end{minipage}

\paragraph{An integer form of the certificate.}
The accompanying verification also computes the same entries by a
separate integer recurrence. Set
\[
 P_{0,n}=\sum_{j=0}^n\binom nj3^{\binom j2}
                    10^{\binom n2-\binom j2},
 \qquad
 P_{r+1,n}=10^{2^r}P_{r,n}^{\,2}-P_{r,n-1}P_{r,n+1}.
\]
Define $d_{0,n}=\binom n2$ and, for $r\geq1$,
$d_{r,n}=2^r\binom n2+r2^{r-1}$.
The identities
\[
 d_{r+1,n}=2d_{r,n}+2^r
          =d_{r,n-1}+d_{r,n+1}
\]
show inductively that $a_{r,n}=P_{r,n}/10^{d_{r,n}}$.
All 49 entries in the triangular calculation agree between the two
recurrences, and all ten verification groups pass.
The terminal rational value and the three preceding values are saved
in full in the verification record.

\paragraph{Consequence for the stronger coefficient conjecture.}
The same source's Conjecture 3.9 asserts positivity of the coefficients
of $\LL^k(f)_n(z)/(1-z)^{2^k-1}$.
At $k=n=6$ its value at $z=3/10$ is negative, because the denominator
$(7/10)^{63}$ is positive. Thus that stronger assertion fails as well.
This is a consequence of the present counterexample, with no additional
solve counted here.

\paragraph{Task 2. Determine the parameter set.}
Characterize those $z\in[0,1]$ for which $(f_n(z))$ is infinitely
log-concave.

\paragraph{Task 3. Determine the first failing iteration.}
For each parameter where infinite log-concavity fails, determine the
least iteration and the first index witnessing failure.

\paragraph{Task 4. Study the normalized coefficient signs.}
Determine when the normalized iterates
$\LL^k(f)_n(z)/(1-z)^{2^k-1}$ have nonnegative coefficients, and explain
the onset of negative coefficients.

\paragraph{Task 5. Find sufficient conditions for binomial transforms.}
Identify useful hypotheses on a nonnegative sequence $(b_j)$ ensuring
infinite log-concavity of $\bigl(\sum_j\binom nj b_j\bigr)_{n\geq0}$.

\begin{thebibliography}{9}
\bibitem{DU}
K. Dilcher and M. Ulas, \emph{Some properties of a class of sparse
polynomials}, Journal of Mathematical Analysis and Applications
\textbf{505} (2022), 125449.
Conjectures 3.8 and 3.9, published p.~10; preprint p.~9,
\href{https://arxiv.org/abs/2008.01480}{arXiv:2008.01480}.
\href{https://doi.org/10.1016/j.jmaa.2021.125449}{Published article}.
\end{thebibliography}
\end{document}
